\subsection{Linear preparation for fixed-width scores}
\label{sec:radix}
The comparison-based preparation can be replaced when scores have the
artifact's finite IEEE-754 binary64 representation. Counting radix sort is
standard~\cite{demaine2011radix}; order-preserving floating-point bit
transforms also predate this work~\cite{herf2001radix}. We combine these
existing techniques with the exact support-event representation and a linear
interval sweep. This is a VUS-specific composition, not a new sorting method.

After the public Float64 conversion, normalize both signed zeros to the
positive-zero bit pattern. Let $u$ be the resulting unsigned 64-bit word and
$h=2^{63}$. Form an ascending key by complementing every bit when the sign
bit is set, and by XOR with $h$ otherwise; complement this key for descending
score order. Eight stable counting passes, from the least to the most
significant byte, sort position indices using 256 buckets per pass. No
comparison-sort library dispatch is assumed.

Scan the resulting permutation by complete global tie groups. A group
containing a support position becomes one retained event; its starting and
ending ranks give $c_j^>$ and $c_j^\ge$. All support positions in that group
receive its event index. A group containing no support position advances the
global rank without creating an event. Consequently off-support points tied
with a support score enter together with that score, rather than being
mistakenly charged to the preceding negative-only rectangle.

For geometry, scan the temporally ordered support positions with a pointer
to the first original interval whose start exceeds the current position.
Inspect the two immediately preceding and two immediately following
intervals, when present, and retain the two smallest distances. Advancing
the pointer over all support positions takes $O(Q+K)$ time, where $K$ is the
number of original anomaly intervals.

\begin{proposition}[Fixed-width score preparation]
\label{prop:radix}
For finite binary64 score keys, assume unit-cost word operations, addresses
and counts that fit a machine word, and unit-cost metric arithmetic including
square roots. Preparation takes $O(n+Q+K)$ time. Exhaustive evaluation of
either scalar PR or the ROC--PR pair takes $O(n+(L+1)Q)$ time and preserves
the target integrals in exact metric arithmetic. The paired implementation
uses $O(n+Q)$ auxiliary words; the current PR-only implementation uses
$O(n+Q+L)$ because it retains one scalar per width.
\end{proposition}
\begin{proof}
For nonnegative finite binary64 values the unsigned bit patterns increase
with numerical value; for negative values their magnitude order is reversed.
The sign-dependent transform therefore preserves numerical order, and the
final complement reverses it. Canonicalizing zero makes key equality
equivalent to numerical equality on the validated finite domain. Stability
of each counting pass preserves the previously sorted lower bytes, so eight
passes give the complete descending order in $O(8(n+256))$ time. The tie scan
then obtains the exact global counts and support groups in $O(n+Q)$ time.

Original intervals are disjoint and ordered. Among intervals starting at or
before a position, distance cannot decrease when moving farther left; among
those starting after it, distance cannot decrease when moving farther right.
Thus two neighbors from each side contain two nearest distinct intervals,
including the possible interval containing the position. The monotone
pointer advances at most $K$ times. The two-distance lemma gives the same
summed-and-clipped weights at every width.

The remaining width kernel is unchanged: it processes $Q$ support positions,
$G$ score events and at most $K$ intervals per width. Since $G,K\le Q$, its
total cost is $O((L+1)Q)$. The exact event formulas preserve the omitted
negative thresholds and the terminal ROC transition. Input-sized keys,
permutation, scratch and position-map arrays require $O(n)$ words; support
and interval state requires $O(Q)$ more. The paired kernel streams the width
sum, whereas the PR-only completion interface stores $L+1$ scalar values.
\end{proof}

This bound uses a fixed 64-bit key universe; it is not a comparison-model
claim for arbitrary real scores or a bit-complexity bound for growing score
precision. Float64 conversion defines the threshold ties, including those
created when distinct large integers convert to the same score. The key and
count construction is discrete; subsequent floating-point accumulation has
the same numerical, rather than symbolic, exactness limitations as the
preceding evaluators.
